% =============================================================================
\section{Introduction}
\label{sec:intro}
% =============================================================================
% План: (1) мотивация -> (2) почему существующие платформы не подходят ->
% (3) наш подход -> (4) вклад (bullet list) -> (5) структура статьи

% --- (1) Мотивация ---------------------------------------------------------
Accurate, low-latency ego-motion estimation and mapping is a prerequisite for
autonomous navigation of micro aerial vehicles, mobile robots, and
head-mounted AR/VR devices~\cite{cadena2016past}.
Modern feature-based and direct SLAM systems such as
ORB-SLAM3~\cite{campos2021orbslam3} and VINS-Mono~\cite{qin2018vins}
achieve centimeter-level accuracy, but \TODO{quantify: power / latency on
embedded CPU (e.g., \NUM~W on Jetson-class SoC)}.

% --- (2) Проблема ------------------------------------------------------------
\TODO{Explain the gap: power budget of target platforms (e.g., nano-drones
$<$\NUM~mW for compute), memory bandwidth limits, irregular sparse workloads
in the back-end that map poorly to GPUs.}
Prior accelerators~\cite{suleiman2019navion,li2019cnn_slam_asic}
\TODO{what they did and what they miss}.

% --- (3) Подход --------------------------------------------------------------
To address these challenges, we propose \sysname, \TODO{one-paragraph
description of the architecture and the central insight}.

% --- (4) Вклад -------------------------------------------------------------
The main contributions of this work are:
\begin{itemize}
  \item \TODO{Contribution 1: algorithm--hardware co-design, e.g.,
        hardware-friendly reformulation of \ldots};
  \item \TODO{Contribution 2: micro-architecture, e.g., streaming front-end
        with \NUM~KB line buffers instead of full-frame storage};
  \item \TODO{Contribution 3: silicon implementation and measurement in
        \NUM-nm CMOS};
  \item \TODO{Contribution 4: comprehensive evaluation against CPU, GPU,
        FPGA and ASIC baselines on public datasets}.
\end{itemize}

% --- (5) Структура -----------------------------------------------------------
The remainder of the paper is organized as follows.
\Cref{sec:related} reviews related work.
\Cref{sec:background} summarizes the SLAM pipeline and profiles its
computational bottlenecks.
\Cref{sec:architecture} presents the proposed architecture, and
\cref{sec:implementation} describes its physical implementation.
\Cref{sec:evaluation} reports experimental results,
\cref{sec:discussion} discusses limitations, and
\cref{sec:conclusion} concludes the paper.
